% ==============================================================================
\section{Computational Experiments}
\label{sec:experiments}
% ==============================================================================

This section compares the IGVND heuristic of \S\ref{sec:heuristic} against
the exact MILP of \S\ref{sec:milp} on a controlled benchmark.  The MILP is
solved with Gurobi and serves as the reference: where it converges it
provides a proven optimum; where it does not, its 60\,s incumbent is the
best exact-solver result obtainable under the same budget.

% ------------------------------------------------------------------------------
\subsection{Benchmark design}
\label{sec:benchmark}
% ------------------------------------------------------------------------------

The benchmark isolates three axes of difficulty: the blocking topology, the
fleet size $R$, and the time-window slack.  All instances use $P = 5$
positions.  Five topology families span the spectrum from independent to
constrained hangars: \emph{None} ($\calA = \emptyset$), \emph{Chain}
($\calA = \{(p_i, p_{i+1})\}$), \emph{Hub} (one aisle position blocks all
others), \emph{Triangle} (a transitive triangle on the rear trio, the layout
of the facility that motivated this study), and \emph{Two rows} (front and
rear halves paired position-by-position).  \emph{None}, \emph{Chain} and
\emph{Hub} are evaluated at the baseline ($R = 10$, tight windows) to isolate
the effect of blocking structure.  \emph{Two rows} and \emph{Triangle} are
additionally crossed over four fleet sizes ($R \in \{5, 10, 20, 30\}$) and
three slack levels (loose, medium, tight), so the effects of size and of
time-window tightness can each be read independently.  This gives
\textbf{27 configurations}, each replicated with 10 independently seeded
instances, for \textbf{270 instances} in total (Table~\ref{tab:benchmark}).
Processing times, earliest starts and due dates are sampled on an integer-day
grid; the tow time is fixed at $\varepsilon = 0.5$ days.

\begin{table}[htbp]
  \centering
  \caption{Benchmark specification: 27 configurations $\times$ 10 seeds
           $= 270$ instances, all with $P = 5$ positions.}
  \label{tab:benchmark}
  \begin{tabular}{lccc}
    \toprule
    Topology & $R$ & Slack & Configs \\
    \midrule
    None / Chain / Hub & 10 & tight & 3 (one each) \\
    Two rows & $\{5, 10, 20, 30\}$ & $\{$loose, medium, tight$\}$ & 12 \\
    Triangle & $\{5, 10, 20, 30\}$ & $\{$loose, medium, tight$\}$ & 12 \\
    \midrule
    \multicolumn{3}{l}{Total} & 27 configs $\times$ 10 seeds $=$ 270 \\
    \bottomrule
  \end{tabular}
\end{table}

% ------------------------------------------------------------------------------
\subsection{Experimental setup}
\label{sec:exp_setup}
% ------------------------------------------------------------------------------

Both methods are run on every instance under three weight profiles that
each isolate one objective term: makespan-priority
$w^{\mathrm{MK}} = (100, 1, 1)$, delay-priority
$w^{\mathrm{DLY}} = (1, 100, 1)$, and movement-priority
$w^{\mathrm{MOV}} = (1, 1, 100)$, for $(W^M, W^D, W^S)$.  Each run is given
a strictly enforced wall-clock budget of 60\,s.  We report the relative
objective gap
\[
  g = \frac{z_{\mathrm{MILP}} - z_{\mathrm{IGVND}}}{z_{\mathrm{MILP}}},
\]
so that $g > 0$ means IGVND attains the lower (better) objective on this
minimisation problem.  Because the relative gap is distorted by small
denominators when the optimal delay is near zero, we also report the
undistorted per-component differences (IGVND $-$ MILP) for makespan, delay
and movement count.  Both quantities are aggregated as the mean over the
10 seeds of each configuration.  All IGVND solutions are validated by the
independent feasibility checker before being scored; across the full
$270 \times 3 = 810$ heuristic runs there were no feasibility failures.

% ------------------------------------------------------------------------------
\subsection{MILP convergence}
\label{sec:milp_conv}
% ------------------------------------------------------------------------------

Table~\ref{tab:milp_conv} reports the mean optimality gap Gurobi still
reports at the end of the budget, per configuration and profile.  The pattern
frames the entire comparison.  The MILP is reliable as a \emph{proven optimum}
(mean gap $0.0\%$) only on the small and loosely constrained cases: the
five-aircraft instances and the no-blocking baseline close on all profiles in
well under a second, as does the loosely-constrained $R{=}10$ \emph{Two rows}
configuration; on the loose $R{=}10$ \emph{Triangle} only movement-priority
still closes on every seed.  As
soon as the fleet grows ($R \ge 20$) or the profile activates the manoeuvre
machinery (\emph{delay-} and \emph{movement-priority} couple the timing and
classification binaries far more tightly than makespan-priority), the gap blows
up: it exceeds $60\%$ at $R{=}20$, climbs to around $78$--$84\%$ at $R{=}30$, and reaches
essentially $100\%$ under the delay-priority profile, where Gurobi returns a
feasible incumbent without closing the bound at all.  Makespan-priority is
consistently the easiest profile for the exact solver (smallest gaps),
delay-priority the hardest.

Even where it cannot close the bound, the exact solver returned a feasible
incumbent within the budget on every run of the benchmark.  On every
configuration with a non-zero gap, a positive $g$ in
Table~\ref{tab:gap_profile} therefore does not mean IGVND beats a proven
optimum; it means IGVND finds a better feasible solution than the exact
solver does in the same time.

\begin{table}[htbp]
  \centering
  \caption{Mean optimality gap reported by Gurobi within the 60\,s budget, per configuration and weight profile (in \%, over 10 seeds). $0.0$ means a proven optimum on every seed; a larger value means Gurobi returns a feasible incumbent without closing the bound. A superscript $(k)$ marks the $k$ seeds on which Gurobi \emph{ran out of memory and returned no solution at all} -- the MILP could not solve the instance; the reported value averages the remaining seeds. \textsc{oom} marks a configuration the MILP could not solve on any seed.}
  \label{tab:milp_conv}
  \begin{tabular}{lll rrr}
    \toprule
    Type & $R$ & Slack & $w^{\mathrm{MK}}$ & $w^{\mathrm{DLY}}$ & $w^{\mathrm{MOV}}$ \\
    \midrule
    None & 10 & tight & 0.0 & 0.0 & 0.0 \\
    Chain & 10 & tight & 28.6 & 99.6 & 82.0 \\
    Hub & 10 & tight & 2.7 & 84.2 & 31.9 \\
    Two rows & 5 & loose & 0.0 & 0.0 & 0.0 \\
     &  & medium & 0.0 & 0.0 & 0.0 \\
     &  & tight & 0.0 & 0.0 & 0.0 \\
     & 10 & loose & 0.0 & 0.0 & 0.0 \\
     &  & medium & 0.0 & 99.2 & 25.4 \\
     &  & tight & 0.0 & 99.6 & 51.3 \\
     & 20 & loose & 60.4 & 99.9 & 90.0 \\
     &  & medium & 61.2 & 99.9 & 92.6 \\
     &  & tight & 61.6 & 99.9 & 93.7 \\
     & 30 & loose & 76.2 & 100.0 & 96.6 \\
     &  & medium & 74.8 & 100.0 & 97.0 \\
     &  & tight & 75.1$^{(1)}$ & 100.0 & 97.4 \\
    Triangle & 5 & loose & 0.0 & 0.0 & 0.0 \\
     &  & medium & 0.0 & 0.0 & 0.0 \\
     &  & tight & 0.0 & 0.0 & 0.0 \\
     & 10 & loose & 0.0 & 0.0 & 0.0 \\
     &  & medium & 0.0 & 99.3 & 42.5 \\
     &  & tight & 0.0 & 99.6 & 58.2 \\
     & 20 & loose & 65.2 & 99.9 & 91.0 \\
     &  & medium & 65.3 & 99.9 & 93.3 \\
     &  & tight & 67.0 & 99.9 & 94.3 \\
     & 30 & loose & 78.7$^{(1)}$ & 100.0 & 97.2 \\
     &  & medium & 80.2 & 100.0 & 97.7 \\
     &  & tight & 80.7 & 100.0 & 97.9 \\
    \bottomrule
  \end{tabular}
\end{table}


% ------------------------------------------------------------------------------
\subsection{Results}
\label{sec:results}
% ------------------------------------------------------------------------------

Table~\ref{tab:gap_profile} reports the relative gap and
Table~\ref{tab:components} the per-component differences.  We read them
together, grouping the configurations by what the MILP can certify.

\paragraph{Instances the MILP solves to optimality.}
On the configurations the MILP solves to optimality, IGVND is exact or
concedes only a few objective units.  It is \emph{exact} on the
no-blocking baseline and on the loose and medium five-aircraft
configurations --- $g = 0.00\%$ on every seed under all three profiles ---
and near-exact on the tight ones, where it concedes a single extra
manoeuvre or a couple of delay units on at most three of the ten seeds (mean gaps
between $-0.6\%$ and $-2.0\%$ under the delay- and movement-priority
profiles, $0.0\%$ under makespan-priority).  On the converged $R{=}10$
cells --- \emph{Two rows}-loose under every profile, \emph{Triangle}-loose
under movement-priority --- the concessions are equally small in absolute
terms (mean $\Delta m \le +1.6$ and $\Delta D \le +2.3$ in
Table~\ref{tab:components}), although the relative gap column can look
large because of the small denominators discussed next.  Under
makespan-priority, where the remaining $R{=}10$ topologies are only
partially converged (Table~\ref{tab:milp_conv}), IGVND tracks the MILP
reference within $-3.5\%$ everywhere.

\paragraph{Delay-priority outliers.}
The two loose $R{=}10$ rows under delay-priority are the extremes of
Table~\ref{tab:gap_profile}: $-14\%$ on \emph{Two rows} and $-129\%$ on
\emph{Triangle}.  Neither reflects a large physical difference.  On these
configurations the optimal delay is zero or near zero, so
$z_{\mathrm{MILP}}$ reduces to the lightly weighted makespan and movement
terms, and a difference of a few delay units --- weighted $100$ ---
swamps that baseline.  On \emph{Triangle}, where the MILP proves
optimality on nine of the ten seeds, IGVND concedes at most three delay
units on any seed (mean $\Delta D = +1.0$ in
Table~\ref{tab:components}), yet single seeds register relative gaps of
$-280\%$ and $-397\%$; on \emph{Two rows} the mean $\Delta D$ is $+0.1$
delay units and IGVND matches the certified optimum exactly on three
seeds.  The per-component columns, not the relative gap, are the
meaningful scale on these rows.

\paragraph{Instances the MILP cannot close within the budget.}
On the configurations the MILP cannot close in 60\,s the comparison
splits by scale.  At $R{=}10$ the timed-out incumbents remain
competitive: IGVND beats them clearly under delay-priority on the
blocking topologies (\emph{Chain} $+10.6\%$, winning all ten seeds;
\emph{Hub} $+4.5\%$), while elsewhere the incumbent edges the heuristic
by a few objective units (the movement-priority column is discussed
below).  From $R{=}20$ the picture is uniform: IGVND returns the better
feasible solution on \emph{every} configuration and profile, and the
advantage grows monotonically with fleet size --- relative gaps of
roughly $+13\%$ to $+32\%$ at $R{=}20$ and $+28\%$ to $+56\%$ at
$R{=}30$, across both the \emph{Two rows} and \emph{Triangle} families
and all three slack levels, peaking under delay-priority ($+55.5\%$ on
\emph{Triangle}-loose-$R{=}30$, all ten seeds won).  The per-component
differences confirm a large absolute improvement and not a denominator
effect: mean $\Delta D$ from roughly $-80$ to $-2\,500$ delay units and
$\Delta m$ from tens to over $250$ makespan units, negative throughout.
The mechanism is the architecture of
\S\ref{sec:igvnd_principle}: by searching over the compact combinatorial
state and pricing it with a fast decoder, IGVND evaluates orders of magnitude
more candidate schedules per second than the branch-and-bound search can
while it is still struggling to tighten its bound.

\begin{table}[htbp]
  \centering
  \caption{Relative objective gap of IGVND against the MILP baseline, $g = (z_{\mathrm{MILP}} - z_{\mathrm{IGVND}})/z_{\mathrm{MILP}}$ (mean over 10 seeds, in \%). $g>0$ means IGVND attains the \emph{lower} objective. All runs share a 60\,s budget. Table~\ref{tab:milp_conv} reports the MILP's optimality gap per configuration; where it is large the MILP reference is a timed-out incumbent, not a proven optimum.}
  \label{tab:gap_profile}
  \begin{tabular}{lll rrr r}
    \toprule
    Type & $R$ & Slack & $w^{\mathrm{MK}}$ & $w^{\mathrm{DLY}}$ & $w^{\mathrm{MOV}}$ & All \\
    \midrule
    None & 10 & tight & +0.0 & +0.0 & +0.0 & +0.0 \\
    Chain & 5 & loose & +0.0 & +0.0 & +0.0 & +0.0 \\
     &  & medium & +0.0 & +0.0 & +0.0 & +0.0 \\
     &  & tight & -0.0 & -25.2 & +0.0 & -8.4 \\
     & 10 & loose & -4.4 & -111.4 & -4.4 & -40.1 \\
     &  & medium & -4.1 & +0.8 & -2.0 & -1.8 \\
     &  & tight & -1.6 & +13.9 & +2.0 & +4.8 \\
     & 20 & loose & +15.5 & +32.1 & +15.2 & +20.9 \\
     &  & medium & +17.2 & +25.0 & +19.9 & +20.7 \\
     &  & tight & +13.8 & +18.5 & +16.0 & +16.1 \\
     & 30 & loose & +30.9 & +43.2 & +37.1 & +37.0 \\
     &  & medium & +32.1 & +36.5 & +30.1 & +32.9 \\
     &  & tight & +31.5 & +34.8 & +31.7 & +32.6 \\
    Hub & 5 & loose & +0.0 & +0.0 & +0.0 & +0.0 \\
     &  & medium & +0.0 & +0.0 & +0.0 & +0.0 \\
     &  & tight & +0.0 & -2.9 & +0.0 & -1.0 \\
     & 10 & loose & -2.5 & -148.0 & +0.0 & -50.2 \\
     &  & medium & -2.5 & -1.4 & +0.0 & -1.3 \\
     &  & tight & -2.0 & +6.5 & +2.8 & +2.4 \\
     & 20 & loose & +13.4 & +22.2 & +25.4 & +20.3 \\
     &  & medium & +12.5 & +18.6 & +20.8 & +17.3 \\
     &  & tight & +13.2 & +15.1 & +17.5 & +15.3 \\
     & 30 & loose & +14.5 & +23.3 & +22.9 & +20.2 \\
     &  & medium & +15.3 & +19.4 & +18.1 & +17.6 \\
     &  & tight & +15.2 & +18.4 & +17.6 & +17.0 \\
    Two rows & 5 & loose & +0.0 & +0.0 & +0.0 & +0.0 \\
     &  & medium & +0.0 & +0.0 & +0.0 & +0.0 \\
     &  & tight & +0.0 & +0.0 & +0.0 & +0.0 \\
     & 10 & loose & -0.2 & -10.6 & -0.1 & -3.6 \\
     &  & medium & -0.2 & +0.2 & +0.0 & -0.0 \\
     &  & tight & -0.2 & +2.0 & +0.8 & +0.9 \\
     & 20 & loose & +13.9 & +23.0 & +19.1 & +18.7 \\
     &  & medium & +17.1 & +23.5 & +18.8 & +19.8 \\
     &  & tight & +19.8 & +18.7 & +17.0 & +18.5 \\
     & 30 & loose & +30.3 & +43.5 & +33.1 & +35.6 \\
     &  & medium & +29.8 & +36.0 & +36.5 & +34.1 \\
     &  & tight & +30.1 & +32.4 & +34.6 & +32.3 \\
    \bottomrule
  \end{tabular}
\end{table}


\begin{table}[htbp]
  \centering
  \footnotesize
  \caption{Per-component mean difference (IGVND $-$ MILP) over 10 seeds for each weight profile: makespan $\Delta m$, total delay $\Delta D$, movement count $\Delta n$. Negative values favour IGVND. Read alongside Table~\ref{tab:gap_profile}, whose relative gap is inflated by small denominators when the optimal delay is near zero.}
  \label{tab:components}
  \begin{tabular}{lll *{9}{r}}
    \toprule
    & & & \multicolumn{3}{c}{$w^{\mathrm{MK}}$} & \multicolumn{3}{c}{$w^{\mathrm{DLY}}$} & \multicolumn{3}{c}{$w^{\mathrm{MOV}}$} \\
    \cmidrule(lr){4-6}\cmidrule(lr){7-9}\cmidrule(lr){10-12}
    Type & $R$ & Slack & $\Delta m$ & $\Delta D$ & $\Delta n$ & $\Delta m$ & $\Delta D$ & $\Delta n$ & $\Delta m$ & $\Delta D$ & $\Delta n$ \\
    \midrule
    None & 10 & tight & +0.0 & +0.0 & +0.0 & +0.0 & +0.0 & +0.0 & +0.0 & +0.0 & +0.0 \\
    Chain & 5 & loose & +0.0 & +0.0 & +0.0 & +0.0 & +0.0 & +0.0 & +0.0 & +0.0 & +0.0 \\
     &  & medium & +0.0 & +0.0 & +0.0 & +0.0 & +0.0 & +0.0 & +0.0 & +0.0 & +0.0 \\
     &  & tight & +0.0 & +0.8 & -0.4 & +0.1 & +0.1 & +0.0 & +0.0 & +0.0 & +0.0 \\
     & 10 & loose & +2.8 & -2.8 & +1.8 & +1.9 & +1.9 & +2.6 & +0.6 & +3.0 & +0.0 \\
     &  & medium & +2.5 & +1.2 & +3.4 & +0.4 & -0.6 & +1.4 & +0.5 & +2.5 & +0.0 \\
     &  & tight & +1.1 & -13.2 & +0.6 & -6.1 & -18.9 & -1.2 & -1.0 & -4.5 & +0.0 \\
     & 20 & loose & -25.3 & -189.2 & +17.4 & -31.5 & -201.9 & +18.8 & -25.8 & -97.6 & +0.0 \\
     &  & medium & -27.8 & -175.9 & +16.6 & -35.2 & -202.4 & +18.8 & -41.9 & -174.1 & +0.0 \\
     &  & tight & -22.3 & -120.7 & +15.0 & -32.7 & -172.2 & +18.0 & -35.0 & -159.3 & +0.0 \\
     & 30 & loose & -104.5 & -1117.7 & +66.0 & -106.3 & -1287.8 & +60.4 & -106.5 & -1143.8 & +0.0 \\
     &  & medium & -108.8 & -1264.8 & +64.0 & -102.0 & -1195.3 & +57.2 & -113.3 & -973.8 & +0.0 \\
     &  & tight & -108.4 & -1274.8 & +70.0 & -103.9 & -1255.8 & +58.8 & -107.3 & -1176.8 & +0.0 \\
    Hub & 5 & loose & +0.0 & +0.0 & +0.0 & +0.0 & +0.0 & +0.0 & +0.0 & +0.0 & +0.0 \\
     &  & medium & +0.0 & +0.0 & +0.0 & +0.0 & +0.0 & +0.0 & +0.0 & +0.0 & +0.0 \\
     &  & tight & +0.0 & +0.0 & +0.0 & +0.0 & +0.1 & -0.2 & +0.0 & +0.0 & +0.0 \\
     & 10 & loose & +1.6 & -3.6 & -6.8 & -0.6 & +1.9 & -3.8 & -0.1 & +0.1 & +0.0 \\
     &  & medium & +1.6 & +0.3 & -7.0 & -1.9 & +0.8 & -1.2 & +0.0 & +0.0 & +0.0 \\
     &  & tight & +1.4 & -2.8 & -7.6 & -5.3 & -8.7 & -0.6 & -2.8 & -3.5 & +0.0 \\
     & 20 & loose & -18.9 & -105.5 & -16.8 & -22.6 & -105.5 & -23.4 & -24.6 & -145.6 & +0.0 \\
     &  & medium & -17.6 & -114.0 & -14.2 & -24.8 & -123.3 & -14.4 & -25.0 & -152.2 & +0.0 \\
     &  & tight & -19.3 & -96.3 & -16.6 & -19.1 & -124.8 & -8.6 & -21.9 & -149.2 & +0.0 \\
     & 30 & loose & -38.3 & -330.4 & +16.4 & -39.9 & -453.5 & +8.0 & -53.5 & -461.8 & +0.0 \\
     &  & medium & -40.4 & -435.8 & +14.8 & -46.3 & -451.6 & +11.4 & -48.1 & -429.1 & +0.0 \\
     &  & tight & -41.0 & -376.6 & +21.6 & -45.3 & -481.1 & +6.0 & -42.8 & -460.6 & +0.0 \\
    Two rows & 5 & loose & +0.0 & +0.0 & +0.0 & +0.0 & +0.0 & +0.0 & +0.0 & +0.0 & +0.0 \\
     &  & medium & +0.0 & +0.0 & +0.0 & +0.0 & +0.0 & +0.0 & +0.0 & +0.0 & +0.0 \\
     &  & tight & +0.0 & +0.0 & +0.0 & +0.0 & +0.0 & +0.0 & +0.0 & +0.0 & +0.0 \\
     & 10 & loose & +0.1 & -0.1 & -0.6 & +0.5 & +0.1 & -1.2 & +0.3 & -0.2 & +0.0 \\
     &  & medium & +0.1 & -0.1 & -0.4 & -1.8 & -0.1 & -0.4 & +0.1 & -0.1 & +0.0 \\
     &  & tight & +0.1 & +0.5 & -0.6 & -4.5 & -2.0 & -1.8 & -0.7 & -0.6 & +0.0 \\
     & 20 & loose & -20.2 & -141.4 & -2.6 & -23.1 & -104.2 & -4.6 & -21.4 & -95.8 & +0.0 \\
     &  & medium & -27.1 & -154.6 & -0.6 & -36.5 & -178.7 & -1.4 & -27.5 & -135.8 & +0.0 \\
     &  & tight & -31.9 & -279.6 & -0.6 & -33.0 & -162.2 & -0.6 & -24.2 & -144.0 & +0.0 \\
     & 30 & loose & -101.9 & -1097.0 & +1.8 & -116.0 & -1162.7 & +2.4 & -88.2 & -771.5 & +0.0 \\
     &  & medium & -106.6 & -993.1 & -1.0 & -113.2 & -1123.2 & +0.2 & -133.2 & -1293.2 & +0.0 \\
     &  & tight & -109.7 & -1115.5 & +7.4 & -121.8 & -1089.3 & -6.2 & -127.2 & -1273.3 & +0.0 \\
    \bottomrule
  \end{tabular}
\end{table}


\paragraph{Movement-priority profile.}
The movement-priority profile is the hardest for a heuristic.  Both methods
almost always reach a zero-manoeuvre schedule ($\Delta n = 0$ throughout), so
the heavily-weighted movement term contributes nothing and the objective
collapses to makespan plus delay; the winner is then simply whichever packs
the zero-movement schedule tightest.  Because that denominator is not small
(it ranges from $\approx 36$ at $R{=}5$ to several thousand at $R{=}30$), the
relative gaps here are genuine and not a small-denominator artefact --- unlike
the delay-priority outliers of the previous paragraph.  IGVND matches the MILP
\emph{exactly} on every loose and medium five-aircraft configuration
($g = 0.00\%$, all seeds)
and dominates at $R{=}20$ and $R{=}30$ ($+14\%$ to $+54\%$, with large negative
$\Delta m$ and $\Delta D$ --- hundreds of objective units in absolute terms).
Its weakest cells are the $R{=}10$ blocking topologies.  On the loose
\emph{Triangle}, which the MILP solves to proven optimality, IGVND
concedes a few units of makespan or delay on nine of the ten seeds
(mean $g = -5.7\%$; mean $\Delta m = +1.6$, $\Delta D = +2.3$); on the
tight \emph{Chain} it concedes a mean $+32$ delay units to an incumbent
the solver itself certifies only to a $79\%$ optimality gap
($g = -18.9\%$; \emph{Hub}: $-1.9\%$ against a $51\%$ gap).  Two
ingredients keep the converged cells at the few-units scale: a
nest-stretch candidate that reproduces the optimum's stay-stretching
mechanism over the actual blocking digraph (\S\ref{sec:igvnd_safety}), and
the random half of the destruction step, which keeps the walk from anchoring
in one basin.  The larger concession on \emph{Chain} is concentrated in
the delay term and is the clearest remaining headroom for the heuristic
at this scale.

% ------------------------------------------------------------------------------
\subsection{Discussion}
\label{sec:discussion}
% ------------------------------------------------------------------------------

The two methods are complementary rather than competing.  The MILP is the
instrument of record on instances it can close: it certifies optimality on
the small and loosely constrained cases, and it is the yardstick that lets
us \emph{measure} the heuristic's quality there --- where IGVND is shown to
be exact or near-exact.  At $R{=}10$ its timed-out incumbents remain
competitive on the blocking topologies.  IGVND is the instrument of choice
at operational scale: on every configuration at $R{=}20$ and beyond, the
heuristic delivers a markedly better feasible schedule in the same 60\,s,
with a guaranteed-feasible zero-movement floor and a checker-validated
manoeuvre-aware polish.  Read jointly, the MILP validates the heuristic on
the tractable end of the benchmark and the heuristic extends usable
solution quality to the intractable end --- which is the regime an
operator actually faces.
